\documentclass[12pt, a4paper]{article}
\usepackage[T1]{fontenc}
\usepackage[utf8]{inputenc}
\usepackage[polish]{babel}
\usepackage[margin=2.5cm]{geometry}
\usepackage{array}
\usepackage{tabularx}
\usepackage{setspace}
\usepackage{verbatim}
\usepackage[hidelinks]{hyperref}
\usepackage{tocloft}
\usepackage{float}
\usepackage{enumitem}
\usepackage{listings}

\renewcommand{\cftsecleader}{\cftdotfill{\cftdotsep}}

\onehalfspacing

\lstset{
    basicstyle=\ttfamily\small,
    breaklines=true,
    frame=single,
    captionpos=b,
    extendedchars=true,
    keepspaces=true,
    columns=fixed,
    commentstyle=\normalfont\ttfamily,
    showstringspaces=false,
    literate={ą}{{\k{a}}}1 {ć}{{\'c}}1 {ę}{{\k{e}}}1 {ł}{{\l{}}}1 {ń}{{\'n}}1 {ó}{{\'o}}1 {ś}{{\'s}}1 {ź}{{\'z}}1 {ż}{{\.z}}1 {Ą}{{\k{A}}}1 {Ć}{{\'C}}1 {Ę}{{\k{E}}}1 {Ł}{{\L{}}}1 {Ń}{{\'N}}1 {Ó}{{\'O}}1 {Ś}{{\'S}}1 {Ź}{{\'Z}}1 {Ż}{{\.Z}}1
}

\title{\vspace*{3cm} \textbf{Dokumentacja Implementacyjna Java}}
\author{Autorzy: Mateusz Bombola \\ Kajetan Karwacki}
\date{\today}

\begin{document}

\maketitle
\thispagestyle{empty}

\newpage
\tableofcontents
\newpage

\section{Cel dokumentu}
Dokument ten stanowi szczegółową specyfikację implementacyjną interfejsu graficznego (GUI) w języku Java, przeznaczonego do wizualizacji grafów planarnych. Opisuje on architekturę wewnętrzną aplikacji, zastosowane klasy, struktury danych oraz techniki integracji z zewnętrznym silnikiem obliczeniowym napisanym w języku C.

\section{Architektura systemu i podział klas}
Aplikacja opiera się na podejściu obiektowym i wykorzystuje bibliotekę Swing do budowy interfejsu użytkownika. System podzielono na warstwę modelu danych oraz warstwę widoku i sterowania.

\subsection{Struktury modelu danych}
Do przechowywania w pamięci informacji o strukturze grafu zdefiniowano dwie klasy pomocnicze:
\begin{itemize}
    \item \textbf{\texttt{Node}:} Reprezentuje wierzchołek grafu. Przechowuje publiczne pola \texttt{id} (typ \texttt{int}) oraz współrzędne kartezjańskie \texttt{x} i \texttt{y} (typ \texttt{double}).
    \item \textbf{\texttt{Edge}:} Reprezentuje krawędź łączącą dwa wierzchołki. Zawiera nazwę \texttt{name} (\texttt{String}), identyfikatory węzłów \texttt{u} i \texttt{v} (\texttt{int}) oraz wagę połączenia \texttt{weight} (\texttt{double}).
\end{itemize}

\subsection{Logika aplikacji i interfejs graficzny}
Główna mechanika programu zawarta jest w dwóch podstawowych klasach zarządzających oknem i renderowaniem grafiki:
\begin{itemize}
    \item \textbf{\texttt{GraphApp} (dziedziczy po \texttt{JFrame}):} Klasa główna inicjalizująca interfejs użytkownika. Odpowiada za budowę paska menu, panelu bocznego oraz logikę zdarzeń (obsługa przycisków, pól tekstowych). Zarządza procesem komunikacji z plikiem wykonywalnym języka C.
    \item \textbf{\texttt{GraphPanel} (dziedziczy po \texttt{JPanel}):} Klasa silnika renderującego. Nadpisuje metodę \texttt{paintComponent(Graphics g)} w celu rysowania węzłów i krawędzi przy pomocy \texttt{Graphics2D}. Przechowuje struktury danych topologii: węzły w słowniku \texttt{Map<Integer, Node>} (zapewniając szybki dostęp po indeksie) oraz krawędzie w liście \texttt{List<Edge>}.
\end{itemize}

\section{Szczegóły implementacyjne modułu renderującego}

\subsection{Rysowanie elementów i transformacje}
Proces renderowania odbywa się z wykorzystaniem antyaliasingu (\texttt{RenderingHints.KEY\_ANTIALIASING}). 
Aplikacja korzysta z obiektu \texttt{AffineTransform} do płynnego przesuwania układu współrzędnych (translacja za pomocą zmiennych \texttt{px} i \texttt{py}) oraz powiększania/pomniejszania widoku (skalowanie za pomocą zmiennej \texttt{zoom}).

Krawędzie rysowane są z dynamiczną zmianą koloru (gradient) od zielonego do czerwonego, gdzie odcień zależy od wagi krawędzi znormalizowanej w przedziale między minimalną a maksymalną załadowaną wagą.

\subsection{Inteligentna adaptacja skali}
Zaimplementowano mechanizm \texttt{getDynamicScale()}, który rozwiązuje problem mikroskopijnych współrzędnych generowanych przez niektóre algorytmy (np. Tutte Embedding). Funkcja skanuje wyliczone pozycje i w przypadku, gdy największa bezwzględna wartość współrzędnej węzła jest mniejsza niż 15, nakłada na cały układ mnożnik (skalę) wynoszącą 150.0. Zabezpiecza to przed utratą precyzji podczas ostatecznego rzutowania współrzędnych zmiennoprzecinkowych na wartości całkowitoliczbowe (\texttt{int}) przed narysowaniem pikseli na ekranie.

\subsection{Automatyczne centrowanie kamery}
Metoda \texttt{resetView()} wykorzystuje algorytm poszukiwania skrajnych punktów wczytanego grafu w celu wyznaczenia jego matematycznego obrysu oraz środka ciężkości. Po odnalezieniu wartości \texttt{minX}, \texttt{maxX}, \texttt{minY}, \texttt{maxY}, obliczane jest przesunięcie panelu (\texttt{px}, \texttt{py}), które środkuje wyrenderowany układ na ekranie użytkownika niezależnie od wyjściowej skali grafu.

\section{Integracja z silnikiem obliczeniowym C}
Integracja między środowiskiem Java a programem w języku C zrealizowana jest za pomocą wywołań systemowych w obrębie klasy \texttt{GraphApp}.

\subsection{Algorytm wyszukiwania ścieżki}
Aby uniknąć konieczności ręcznego definiowania ścieżki do pliku wykonywalnego przez użytkownika, przed uruchomieniem silnika program skanuje zdefiniowaną tablicę ścieżek relatywnych (\texttt{possiblePaths}). Poszukiwany jest plik \texttt{graph\_layout.exe} w folderze roboczym Javy, dedykowanym folderze \texttt{Algorytm c/pliki\_zrodlowe\_2\_grupy/} oraz odpowiednich folderach nadrzędnych. W przypadku niepowodzenia aktywowany jest plan awaryjny wywołujący graficzne okno \texttt{JFileChooser}.

\subsection{Wywołanie procesu i asynchroniczność}
Aplikacja wykorzystuje klasę \texttt{ProcessBuilder} do złożenia komendy wywołującej. Parametry przekazywane są w postaci listy struktur \texttt{String}, co ułatwia zarządzanie argumentami programu docelowego:
\begin{lstlisting}[language=Java]
List<String> command = new ArrayList<>();
command.add(enginePath);
command.add(inputFile.getAbsolutePath());
command.add("-o");
command.add(outputFile.getAbsolutePath());
if (isTutte) {
    command.add("-a");
}
ProcessBuilder pb = new ProcessBuilder(command);
\end{lstlisting}
Po uruchomieniu procesu Java wywołuje metodę \texttt{p.waitFor()}, zawieszając wykonanie bieżącego wątku do momentu zakończenia obliczeń. Strumień błędów programu docelowego (\texttt{getErrorStream()}) jest stale nasłuchiwany za pomocą klasy \texttt{BufferedReader}, co umożliwia natychmiastowe przechwycenie i przekazanie użytkownikowi potencjalnych komunikatów z silnika C.

\section{Mechanizmy wejścia/wyjścia (I/O)}
Moduł I/O obsługuje ładowanie plików tekstowych ze strukturą krawędzi i wynikowych współrzędnych za pomocą klas \texttt{FileReader} oraz \texttt{BufferedReader}. Zapis nowych współrzędnych realizowany jest poprzez \texttt{PrintWriter}. Dodatkowo aplikacja udostępnia metodę \texttt{exportToPNG()}, która poprzez wyrysowanie zawartości do obiektu \texttt{BufferedImage} generuje graficzny zrzut obszaru roboczego do pliku.

\end{document}